% Chapter 8: Future Work and Recommendations
\section{Technical Enhancements}

\subsection{Advanced Reasoning}

Future work could incorporate more sophisticated reasoning mechanisms, including probabilistic reasoning for handling uncertainty in learner models.

\subsection{Scalability Improvements}

Optimization of the triple store and reasoning engine would enable deployment to larger student populations and more complex ontologies.

\subsection{Integration with AI/ML}

Machine learning techniques could enhance student modeling and recommendation generation, learning from aggregate patterns across learners.

\section{Pedagogical Extensions}

\subsection{Multi-disciplinary Learning Paths}

The framework could be extended to model interdisciplinary connections between CS and other domains.

\subsection{Social Learning Features}

Incorporating collaborative learning mechanics and peer recommendation could enhance the social dimensions of learning.

\subsection{Adaptive Difficulty Calibration}

Dynamic difficulty adjustment based on learner performance and confidence could optimize the challenge level for each student.

\section{Broader Adoption and Deployment}

\subsection{Open Source Release}

Making the framework and ontology available as open source resources could accelerate adoption and community-driven improvements.

\subsection{Standards and Interoperability}

Aligning the ontology with educational standards bodies could improve interoperability with existing LMS platforms and educational data systems.

\section{Long-Term Vision}

The ultimate goal is to establish ontology-based approaches as a standard practice in educational technology, enabling more intelligent, adaptive, and equitable learning systems across all educational domains.
